%قبل از شروع مسئله ترجیحاً چیزی ننویسید. اما اگر خواستید بنویسید، در قالب لیستی (فقرات و شمارش) بنویسید.
%برای اضافه کردن نمره، بعدی از شروع‌مسئله، [x] را اضافه کنید و به جای x  نمره مورد نظر را قرار دهید. برای زیرمسئله‌ها هم به همین صورت است.
%\شروع‌مسئله[10]
\شروع‌مسئله
% صورت اصلی مسئله را در این‌جا بنویسید.
درختی دودویی را در نظر بگیرید که در آن هر گره یا دقیقا دو فرزند دارد یا فرزندی ندارد و به گره‌های آن اعداد ۱ تا
 $n$ 
بدون ترتیب خاصی نسبت داده شده است.
پیش‌ترتیب و میان‌ترتیب این درخت به ما داده شده‌است.
روشی برای بدست آوردن خود درخت ارائه کنید و ثابت کنید با داشتن پیش‌ترتیب و میان‌ترتیب، درخت به صورت یکتا مشخص می‌شود.

\solution{
اگر بتوانیم ریشه درخت و فرزندانش را به طور یکتا مشخص کنیم و پیش‌ترتیب و میان‌ترتیب زیردرخت‌های فرزندانش را هم مشخص کنیم، 
مسئله به کمک استقرای قوی روی تعداد رئوس درخت حل می‌شود.
با توجه به تعریف پیش‌ترتیب، عنصر اول پیش‌ترتیب، همان ریشه درخت است. 
ریشه را در میان‌ترتیب در نظر بگیرید. با توجه به اینکه می‌دانیم هر گره یا دو فرزند دارد یا فرزندی ندارد، دنباله‌ی عناصر قبل از ریشه در صورت وجود، الزاما میان‌ترتیب فرزند چپ ریشه است و همچنین عناصر بعد از ریشه در صورت وجود، الزاما میان‌ترتیب فرزند راست ریشه است.
پس حالا ما تمامی گره‌هایی که باید در زیردرخت سمت راست (یا چپ) ریشه باشند را داریم.
با توجه به اینکه عناصر هر زیردرخت را مشخص کرده‌ایم میان‌ترتیب و پیش‌ترتیب‌ زیردرخت راست و چپ هم به سادگی قابل استخراج هستند 
و حال که این موارد را داریم می‌توانیم به صورت بازگشتی عمل کنیم و زیردرخت‌های راست و چپ را بسازیم و بدین ترتیب مسئله حل می‌شود.
برای اثبات یکتا‌ بودن جواب هم دقیقا به همین ترتیب عمل می‌کنیم و از استقرای قوی استفاده می‌کنیم. 

}



%%%%%%%%%%%%%%%%%%%%%%%%%%%%%%%%%%%%%%%%%%%%%%%%%%%%%
% برخی دستورات مفید برای تمرین‌ها. %
%%%%%%%%%%%%%%%%%%%%%%%%%%%%%%%%%%%%%%%%%%%%%%%%%%%%%

% فرمول‌های نمونه:
%$f(x) = \Theta(x)$
%$T(x, y) \leq \log(x) + 2^y + \binom{2n}{n}$
% {\tt code\_function}


% لیست موارد بدون شماره‌گذاری:
%\شروع{فقرات}
%    \فقره مورد اول 
%    \فقره مورد دوم
%\پایان{فقرات}

% لیست موارد شماره‌گذاری شده:
%\شروع{شمارش}
%    \فقره یک مورد 
%    \فقره یک مورد دیگر 
%    \فقره باز هم یک مورد دیگر
%\پایان{شمارش}

% برای نوشتن یک فرمول چند خطی که یک شماره بخوره. برای حالت‌های دیگر به سایت زیر مراجعه کنید:
% https://www.overleaf.com/learn/latex/Aligning_equations_with_amsmath
%\begin{align} 
%    \begin{split}
%        (x+y)^3 &= (x+y)^2(x+y) \\
%                &= (x^2+2xy+y^2)(x+y) \\
%                &= (x^3+2x^2y+xy^2) + (x^2y+2xy^2+y^3) \\
%                &= x^3+3x^2y+3xy^2+y^3
%    \end{split}                                 
%\end{align}

% نوشتن ماتریس و آرایه دوبعدی
%\begin{align}
%    A = 
%    \begin{bmatrix}
%        A_{11} & A_{21} \\
%        A_{21} & A_{22}
%    \end{bmatrix}
%\end{align}

% اضافه کردن تصویر با شرح. برای اضافه کردن تصویر بدون شرح، فقط دستور درج‌تصویر را استفاده کنید.
%\شروع{شکل}[h!]
%\تنظیم‌ازوسط
%  \درج‌تصویر[پهنا=0.5\پهنای‌متن]{img.jpg}
%\شرح{شرح تصویر}
%\برچسب{تصویر:۱}
%\پایان{شکل}
%\پایان{زیرمسائل}

%کد زیر را ببینید.
%%\lstinputlisting{source_filename.py}
%\begin{latin}
%\begin{lstlisting}
%# Import the modules
%import sys
%import random
%
%ans = True
%
%while ans:
%question = raw_input("Ask the magic 8 ball a question: (press enter to quit) ")
%
%answers = random.randint(1,8)
%
%if question == "":
%sys.exit()
%
%elif answers == 1:
%print "It is certain"
%
%elif answers == 2:
%print "Outlook good"
%
%elif answers == 3:
%print "You may rely on it"
%
%elif answers == 4:
%print "Ask again later"
%
%elif answers == 5:
%print "Concentrate and ask again"
%
%elif answers == 6:
%print "Reply hazy, try again"
%
%elif answers == 7:
%print "My reply is no"
%
%elif answers == 8:
%print "My sources say no"
%\end{lstlisting}
%\end{latin}

